% Part: model-theory
% Chapter: basics
% Section: dlo

\documentclass[../../../include/open-logic-section]{subfiles}

\begin{document}

\olfileid{mod}{bas}{dlo}
\section{الترتيبات الخطية الكثيفة}

\begin{defn}
  \emph{الترتيب الخطي الكثيف بلا طرفين} هو !!a{structure}~$\Struct{M}$
  للـ!!{language} التي تحتوي !!{predicate} واحدًا ذا موضعين~$<$ يحقق
  الجمل الآتية:
  \begin{enumerate}
  \item $\lforall[x][\lnot x < x]$؛
  \item $\lforall[x][\lforall[y][\lforall[z][(x < y \lif (y < z \lif x
    <z ))]]]$؛
  \item $\lforall[x][\lforall[y][(x< y \lor \eq[x][y] \lor y < x)]]$؛
  \item $\lforall[x][\lexists[y][x < y]]$؛
  \item $\lforall[x][\lexists[y][y < x]]$؛
  \item $\lforall[x][\lforall[y][(x < y \lif \lexists[z][(x < z \land
        z < y))]]]$.
 \end{enumerate}
\end{defn}

\begin{thm}\ollabel{thm:cantorQ}
  كل ترتيبين خطيين كثيفين !!{enumerable} بلا طرفين متشاكلان.
\end{thm}

\begin{proof}
  لنأخذ $\Struct{M_1}$ و$\Struct{M_2}$ ترتيبين خطيين كثيفين
  !!{enumerable} بلا طرفين، ولنضع ${<_1} = \Assign{<}{M_1}$ و
  ${<_2} = \Assign{<}{M_2}$. ولتكن $\PIso{I}$ جميع التشاكلات
  الجزئية بينهما. فـ $\PIso{I}$ غير خالية، لأن
  $\emptyset \in \PIso{I}$. نريد إثبات الذهاب والإياب لـ $\PIso{I}$؛
  إذ ينتج منهما $\Struct{M_1} \iso[p] \Struct{M_2}$، ثم يكفي
  الرجوع إلى \olref[pis]{thm:p-isom1}.

  أما الذهاب لـ $\PIso{I}$، فنأخذ $p \in \PIso{I}$ و
  $a \in \Domain{M_1}$. إن كانت $p$ خالية، اخترنا أي
  $b \in \Domain{M_2}$ ووضعنا $q = \{\langle a,b\rangle\}$.
  وإن كان $a$ في مجال $p$، كفانا $q=p$. فبقيت حال $p$
  غير الخالية و$a$ خارج مجالها. نكتب $p(a_i)=b_i$، حيث $i=1$،
  \dots،~$n$، ونرتب المدخلات بحيث
  $a_1 <_1 a_2 <_1 \cdots <_1 a_n$، ولا إخلال بذلك بالعموم.
  ثم نختار $b \in \Domain{M_2}$ بحسب الحالات الآتية:
  \begin{enumerate}
  \item إذا كان $a <_1 a_1$، فليكن $b \in \Domain{M_2}$ بحيث
    $b <_2 b_1$؛
  \item إذا كان $a_n <_1 a$، فليكن $b \in \Domain{M_2}$ بحيث
    $b_n <_2 b$؛
 \item إذا كان $a_i <_1 a <_1 a_{i+1}$ لبعض $i$، فليكن
   $b \in \Domain{M_2}$ بحيث $b_i <_2 b <_2 b_{i+1}$.
  \end{enumerate}
  ولا يعسر وجود $b$ على الوجه المطلوب، فإن $\Struct{M_2}$
  كثيف بلا طرفين. نضع $q = p \cup \{ \langle a, b \rangle \}$،
  فتكون $q \in \PIso{I}$ وتوسّع $p$ كما أردنا. فثبت الذهاب،
  والإياب نظيره. وعليه $\Struct{M_1} \iso[p] \Struct{M_2}$؛
  وتنتج من \olref[pis]{thm:p-isom1} العلاقة
  $\Struct{M_1} \iso \Struct{M_2}$.
\end{proof}

\begin{prob}
  أتم برهان \olref[mod][bas][dlo]{thm:cantorQ} بالتحقق من أن
  $\PIso{I}$ تحقق خاصية الإياب.
\end{prob}

\begin{rem}
  كل ترتيب خطي كثيف $\Struct{S}$، إذا كان !!{enumerable} وبلا
  طرفين، لزم من \olref{thm:cantorQ} أن $\Struct{S} \iso \Struct{Q}$،
  حيث $\Struct{Q} = (\Rat, <)$ هو الترتيب الخطي الكثيف
  !!{enumerable} على الأعداد النسبية~$\Rat$. ولنعد إلى
  الـ!!{structure}~$\Struct{R} = (\Real, <)$ في
  \olref[thm]{remark:R}. فقد وجدنا !!a{enumerable}
  !!{structure}~$\Struct{S}$ تحقق $\Struct{R} \elemequiv \Struct{S}$.
  وهذه $\Struct{S}$ ترتيب خطي كثيف !!a{enumerable} بلا طرفين،
  فهي متشاكلة مع الـ!!{structure}~$\Struct{Q}$، ومن ثم مكافئة
  لها أوليًا. وبالتعدي نحصل على $\Struct{R} \elemequiv \Struct{Q}$.
  ولنا طريق مباشر إلى ذلك: نثبت $\Struct{R} \iso[p] \Struct{Q}$
  بحجة الذهاب والإياب المتقدمة نفسها.
\end{rem}
\end{document}
